
\subsubsection{4.1 Dataset}

\textbf{TruthfulQA} \citep{Lin2022TruthfulQA}. The generation subset containing 817 questions designed to elicit plausible but incorrect answers from language models is used. The benchmark covers 38 categories including health, law, finance, and politics.

The dataset is split 80/20 into training (653 questions) and validation (164 questions) using a fixed seed (42). The training split is used to compute semantic entropy labels and train probes; the validation split is held out for all transfer evaluations.

\begin{table}[htbp]
\centering
\resizebox{\columnwidth}{!}{%
\begin{tabular}{lll}
\toprule
Split & Questions & Usage \\
\midrule
Train & 653 & SE label computation, probe training, aligner fitting \\
Validation & 164 & Transfer evaluation (all reported metrics) \\
\bottomrule
\end{tabular}
}
\end{table}

\subsubsection{4.2 Models}

Three instruction-tuned LLMs from different organizations are tested:

\begin{table*}[htbp]
\centering
\resizebox{\dimexpr 2\columnwidth+\columnsep\relax}{!}{%
\begin{tabular}{lllll}
\toprule
Model & Provider & Parameters & Hidden Dim & Layers \\
\midrule
Llama-3-8B-Instruct & Meta & 8B & 4096 & 32 \\
Mistral-7B-Instruct-v0.2 & Mistral AI & 7B & 4096 & 32 \\
Qwen-2-7B-Instruct & Alibaba & 7B & 3584 & 28 \\
\bottomrule
\end{tabular}
}
\end{table*}

\subsubsection{4.3 Evaluation Metrics}

\textbf{AUROC.} Primary metric for hallucination detection. Measures the probability that a randomly chosen correct answer has lower predicted uncertainty than a randomly chosen incorrect answer.

\textbf{Transfer Gap.} For source model i and target model j:
$\text{gap}_{i\rightarrow j} = \text{AUROC}_{i\rightarrow i} - \text{AUROC}_{i\rightarrow j}$

\textbf{Success Criteria.}
\begin{itemize}
\item Mean gap < 0.05: Evidence for architecture-invariant encoding
\item Max gap < 0.10: No pair shows prohibitive transfer degradation
\end{itemize}

\subsubsection{4.4 Experiment Mode}

The experiments reported here used a proof-of-concept validation mode with random binary labels rather than true semantic entropy labels. This is valid for mechanism validation because the transfer gap measures relative performance degradation, not absolute AUROC. A probe trained on random labels learns some discriminative boundary on hidden states; what matters for transfer validation is whether that boundary transfers across models with minimal degradation. Full evaluation with true SE labels is noted as future work.
